
\documentclass{article}

\usepackage{amsmath,amsthm}

\newtheorem{corollary}{Corollary}
\newcommand{\LRF}{\mathrm{LRF}}
\newcommand{\OPT}{\mathrm{OPT}}

\begin{document}

\begin{corollary}
Let $\lambda(I)$ denote the performance ratio of the maximum cost WSPT schedule for an
instance $I$ of $P||\sum_j w_j C_j(\alpha)$, with $ \alpha \in (0,1]$. That is,
\[
        \lambda(I)
        =
        \max_{\text{WSPT schedule }S\text{ of }I}\frac{\sum_{j\in I} w_j C_j(S(I))(\alpha)}
             {\sum_{j\in I} w_j C_j^*(\alpha)},
\]
where $C_j(S(I))(\alpha)$ and $C_j^*(\alpha)$ are the $\alpha$-completion times of job $j$ in the WSPT
schedule  $S$ and in an optimal schedule, respectively. Then for every instance
$I$ of $P||\sum_j w_j C_j(\alpha)$, there is an instance $I'$ of
$P||\sum_j p_j C_j(\alpha)$ with
\[
        \lambda(I')\geq \lambda(I).
\]
\end{corollary}

\begin{proof}
Let
\[
        r_1>r_2>\cdots>r_d
\]
be the distinct ratios $w_j/p_j$ occurring in instance $I$, and set $r_{d+1}=0$.

Note that
\[
        r_i=\sum_{k=i}^d (r_k-r_{k+1}),
\]
Let $i(j)$ be the index  of the ratio for which
\[
        \frac{w_j}{p_j}=r_{i(j)}.
\]

Then it holds that
\[
        w_j
        =
        r_{i(j)}p_j
        =
        \left(\sum_{k=i(j)}^d (r_k-r_{k+1})\right)p_j .
\]
And thus for any schedule S of $I$,
\[
\begin{aligned}
        \sum_{j\in I} w_j C_j(S(I))(\alpha)
        &=
        \sum_{j\in I}
        \left(
            \sum_{k=i(j)}^d (r_k-r_{k+1})
        \right)p_j C_j(S(I))(\alpha)                                      \\
        &\overset{(1)}{=}
        \sum_{k=1}^d (r_k-r_{k+1})
        \sum_{\substack{j\in I\\ i(j)\leq k}} p_j C_j(S(I))(\alpha) .
\end{aligned}
\]

Next we define
\[
        I_k=\{j\in I: i(j)\leq k\} \qquad \text{for } k=1,\ldots,d.
\]
Thus $I_k$ consists of exactly those jobs whose ratios are among the $k$
largest ratios in $I$. Now regard each $I_k$ as an instance of
$P||\sum_j p_j C_j(\alpha)$ by setting $w_j=p_j$ for every job $j\in I_k$.

Let $C(I_k)(\alpha)$ be the value of a cost maximizing WSPT schedule for $I_k$, and
let $C^*(I_k)(\alpha)$ be its optimal value.
The $\alpha$-completion time of every job $j\in I_k$ is the same in this WSPT schedule for
$I_k$ as in the WSPT schedule for $I$ as we have not altered any $p_j$.

Using the decomposition above for the WSPT schedule and for an optimal
schedule of $I$, we obtain
\[
\begin{aligned}
        \lambda(I)
        &=
        \frac{
        \sum_{k=1}^d (r_k-r_{k+1}) C(I_k)
        }{
        \sum_{k=1}^d
        (r_k-r_{k+1})
        \sum_{j\in I_k} p_j C_j^*
        }                                                    \\
        &\overset{(2)}{\le}
        \frac{
        \sum_{k=1}^d (r_k-r_{k+1}) C(I_k)
        }{
        \sum_{k=1}^d
        (r_k-r_{k+1}) C^*(I_k)
        }                                                    \\
        &\leq
        \max_{1\leq k\leq d}
        \frac{C(I_k)}{C^*(I_k)}
        =
        \max_{1\leq k\leq d}\lambda(I_k).
\end{aligned}
\]
Therefore, choosing an index $k$ maximizing $\lambda(I_k)$ and setting
$I'=I_k$ gives an instance of $P\|\sum_j p_j C_j$ with
$\lambda(I')\geq \lambda(I)$.

It remains to show (1) and (2) in detail.
Regarding (1), we show
\[
        \sum_{j\in I}
        \left(
            \sum_{k=i(j)}^d (r_k-r_{k+1})
        \right)p_j C_j(S(I))(\alpha)
        =
        \sum_{k=1}^d \left((r_k-r_{k+1})
        \sum_{\substack{j\in I\\ i(j)\leq k}} p_j C_j(S(I))(\alpha) \right).
\]

by induction on the number of jobs $n = |I|$.

(possibly add this depending on waht maxi says regarding detail)

Regarding (2):

(possibly add this depending on waht maxi says regarding detail)

\end{proof}

\end{document}